% !TeX spellcheck = hu_HU
% !TeX encoding = UTF-8

\chapter*{Összegzés}
\addcontentsline{toc}{chapter}{Összegzés}

A Linda nyelv, bár egy korai párhuzamosítási paradigmára épül a mai napig alkalmas különböző párhuzamos problémák bemutatására és kifejezetten jó oktatási eszközként tud szolgálni.

A dolgozatban megmutattam a párhuzamos programozás szükségességét, és a modern programfejlesztésben való elhanyagolhatatlan szerepét.
Röviden bemutattam a parallelizációs világ hardver és szoftver komponenseit, jobban hangsúlyozva a szoftveres megoldásokat és azok közötti kapcsolatokat, amelyeknek a tárgyalása a Linda nyelvnek, a dolgozat fő témájának, a bevezetéséhez vezetett.

A Linda nyelv elméleti bemutatását követően, ismertettem az általam készített megvalósítás általánosságoktól elindulva elmélyedve a különböző releváns részletekbe.
Elsőként a használt MPI környezet felhasználást a skálázható párhuzamosításra, amely lehetővé teszi az asztali gépektől a szupergépekig való futtatását a Linda megvalósításnak.
A használt adatbázis is részletekbe menően ismertetve lett, hiszen az alkotja a Linda modell központi részét.

Megállapítottam, hogy a Linda kiválóan használható a párhuzamos programozás oktatására, a könnyű mentális modellje és szintaxisa miatt, amivel könnyen el lehet sajátítani, és amin keresztül a parallelizációs technikák bemutatása kifejezetten jól átadhatóak.
Erről az állításomat pár mintapéldával támasztottam alá, amelyeken keresztül a párhuzamosítási elvek egyszerű bemutatása érhetően megvalósítható.

Egy gyakorlati példára adott megvalósításon keresztül megállapítottam, hogy általános párhuzamosításra is alkalmas.
Ezzel bemutattam, hogy a Linda nyelv rövid és átlátható kóddal képes lépést tartani az ipari gyakorlatban sokszor látható eszközökkel, sőt az egyiket stabilan le is győzte a párhuzamos végrehajtásból származó gyorsulás terén.

